\documentclass[11pt]{article}
\usepackage[margin=1in]{geometry}
\usepackage{amsmath,amssymb,amsthm}
\usepackage{mathtools}
\usepackage{booktabs}
\usepackage{xcolor}
\usepackage[colorlinks=true,linkcolor=blue,urlcolor=blue]{hyperref}
\usepackage{enumitem}

\newcommand{\so}[1]{\mathfrak{so}(#1)}
\newcommand{\SO}[1]{\mathrm{SO}(#1)}
\newcommand{\SU}[1]{\mathrm{SU}(#1)}
\newcommand{\Sp}[1]{\mathrm{Sp}(#1)}
\newcommand{\Tr}{\mathrm{Tr}}
\newcommand{\sh}{s_h}
\newcommand{\ch}{c_h}
\newcommand{\half}{\tfrac12}
\newcommand{\code}[1]{\texttt{#1}}
\newcommand{\impl}[2]{\par\smallskip\noindent\textcolor{blue!60!black}{\small$\blacktriangleright$ \textbf{Implemented in:} \code{#1}\quad\textbf{Verified by:} \code{#2}}\par\smallskip}

\title{\textbf{pyCHM derivations}\\[4pt]
\large First-principles group theory of the Minimal and Next-to-Minimal Composite Higgs}
\author{Generated for the \texttt{pyCHM} library}
\date{\today}

\begin{document}
\maketitle

\begin{abstract}
This note derives, from group theory alone, every quantity that \code{pyCHM} treats as
\emph{derived} (as opposed to model input, or numerically anchored to the independent
\code{pypngb} engine). Each step is tied to the function that implements it and the test that
verifies it in closed form, so the document and the code are one object: if a derivation here is
wrong, a named test fails. The reference physics is Murnane, \emph{The Landscape of Composite
Higgs Models} (arXiv:2606.18364); conventions match the library exactly. Nothing here assumes the
thesis is correct --- the closed forms are recomputed from the generators, and where the thesis is
the only source (the App.~A7 form factors) it is cross-checked against the \code{pypngb}-anchored
eigenvalue route (\S\ref{sec:ff}).
\end{abstract}

\tableofcontents

\section{The CCWZ framework}
A composite Higgs is a set of pseudo-Nambu--Goldstone bosons (pNGBs) of a spontaneously broken
global symmetry $G\to H$. Write the generators of $\mathfrak{g}=\mathrm{Lie}(G)$ as the unbroken
$T^a\in\mathfrak{h}$ and the broken $X^{\hat a}\in\mathfrak{g}/\mathfrak{h}$. The pNGBs
$\pi^{\hat a}(x)$ parameterise the coset through the Goldstone matrix
\begin{equation}
  U(\pi) \;=\; \exp\!\Big(\tfrac{i}{f}\,\pi^{\hat a}X^{\hat a}\Big),
  \label{eq:U}
\end{equation}
with $f$ the decay constant. Fermions are embedded in (generally incomplete) representations $R$
of $G$; the Higgs dependence of every elementary--composite mixing is the matrix element of $U$
in $R$,
\begin{equation}
  \big\langle c \,\big|\, U_R(\pi) \,\big|\, E \big\rangle ,
  \label{eq:dressing}
\end{equation}
between the embedding $E$ of the elementary field and a composite state $c$. The central claim of
the library is that \emph{these matrix elements are the only source of the $\sh$-dependence}, and
that they are computable in closed form from the generators --- never fitted. The rest of this
document carries out that computation for $\SO5/\SO4$ (MCHM) and $\SO6/\SO5$ (NMCHM).

\section{Minimal Composite Higgs: \texorpdfstring{$\SO5/\SO4$}{SO(5)/SO(4)}}

\subsection{Generators and conventions}
$\SO5$ acts on indices $0,\dots,4$. The unbroken $\SO4$ acts on $0,\dots,3$; the coset direction
is index $4$. In the vector (the $\mathbf 5$) the Hermitian generators are
\begin{equation}
  \big(T^{AB}\big)_{CD} \;=\; -\,i\big(\delta_{AC}\delta_{BD}-\delta_{AD}\delta_{BC}\big),
  \qquad 0\le A<B\le 4 .
\end{equation}
The six $T^{ab}$ ($a<b\le 3$) generate $\SO4$; the four broken generators are $X^{\hat a}=T^{\hat a 4}$.
With the Higgs vacuum expectation value along the $\hat a=3$ direction, the only field that survives
in unitary gauge is the physical Higgs $h$, and $U$ becomes a single planar rotation.
\impl{ccwz.gen\_vector, symbolic/core.py}{test\_so4\_unbroken\_generators\_close\_the\_algebra}

\subsection{The vector Goldstone matrix}
Exponentiating the broken generator $X^3=T^{34}$ with angle $\theta=h/f$ gives a rotation in the
$(3,4)$ plane,
\begin{equation}
  U_{\mathbf 5}(\theta)=\exp\!\big(i\theta\,T^{34}\big)
  =\begin{pmatrix} \mathbb 1_3 & & \\ & \ch & \sh \\ & -\sh & \ch \end{pmatrix},
  \qquad \sh\equiv\sin\theta,\ \ \ch\equiv\cos\theta=\sqrt{1-\sh^2}.
  \label{eq:U5}
\end{equation}
The library stores everything in the variable $\sh=\sin(h/f)$ with $\ch=\sqrt{1-\sh^2}$, avoiding
an $\arcsin$ round-trip (machine precision). Eq.~\eqref{eq:U5} is the thesis ``explicit Goldstone
matrix.''
\impl{ccwz.U\_vector, core.U\_vector\_sym}{test\_U\_vector\_matches\_thesis\_goldstone}

\subsection{Lifting to tensor representations}
A rank-$k$ $\SO5$ tensor irrep is carried by the totally symmetric-traceless or totally
antisymmetric part of $\mathbf 5^{\otimes k}$. The Goldstone acts on each index,
\begin{equation}
  \big(U_R\big)_{b a}=\big\langle E_b \,\big|\, U_{\mathbf 5}^{\otimes k}\,E_a\big\rangle
  = \Tr\!\big(E_b^{\dagger}\,U E_a U^{\!\top}\big)\quad(\text{rank }2),
  \label{eq:lift}
\end{equation}
in an orthonormal basis $\{E_a\}$ of the irrep. The first three rungs are the
$\mathbf 5$ (rank 1), the $\mathbf{10}$ (antisymmetric rank 2) and the $\mathbf{14}$
(symmetric-traceless rank 2). The construction is dimension-agnostic, so the same code produces
the $\mathbf{30}$ (rank-3 symmetric) and, with $n=6$, the $\SO6$ tower of \S\ref{sec:nmchm}.
\impl{tensors.tensor\_basis, tensors.U\_rep\_tensor, ccwz.U\_rep}{test\_tensors.py (dimension, unitarity, homomorphism)}

\subsection{\texorpdfstring{$\SO4=\SU2_L\times\SU2_R$}{SO(4)=SU(2)xSU(2)} branching}
The unbroken $\SO4$ on indices $0,\dots,3$ is $\SU2_L\times\SU2_R$ via the 't~Hooft combinations
\begin{equation}
  J^{L/R}_i \;=\; \half\Big(\tfrac12\,\epsilon_{ijk}T^{jk}\;\pm\;T^{0i}\Big),\qquad i=1,2,3 .
\end{equation}
Diagonalising the two Casimirs $C_{L}=\sum_i (J^L_i)^2$, $C_{R}=\sum_i (J^R_i)^2$ on the irrep
basis (they commute) labels each sub-multiplet by $(j_L,j_R)$ with $C=j(j+1)$. The branchings are
\begin{align}
  \mathbf 5 &= \big(\tfrac12,\tfrac12\big)\oplus(0,0),\\
  \mathbf{10} &= \big(\tfrac12,\tfrac12\big)\oplus(1,0)\oplus(0,1),\\
  \mathbf{14} &= (1,1)\oplus\big(\tfrac12,\tfrac12\big)\oplus(0,0).
\end{align}
\impl{decompose.so4\_decompose, decompose.\_su2\_casimirs}{test\_branchings\_match\_thesis}

\subsection{Dressings are matrix elements (derive, don't fit)}
A claimed closed form $f(\theta)$ for an elementary--composite mixing is admissible only if it lies
in the span of the single-basis overlaps $\{\langle E_b|U_R|E\rangle\}$. The library solves for the
exact minimal-norm composite tensor $c$ with $\langle c|U_R|E\rangle=f$ and \emph{raises} if $f$ is
not realisable. This is the from-scratch guarantee: every dressing factor in the model files is
\emph{proven} to be a Goldstone matrix element, with symbolic residual exactly $0$.
\impl{symbolic/derive.py::solve\_composite}{test\_dressings\_are\_matrix\_elements}

\subsection{The $\mathbf{5}$-$\mathbf{5}$-$\mathbf{5}$ coefficients}
Embed the up-component of $q_L$ in the bidoublet, $E_{q_L}=(e_0+e_3)/\sqrt2$, and $t_R$ in the
$\SO4$ singlet $E_{t_R}=e_4$. From \eqref{eq:U5}, $U e_4=\sh\,e_3+\ch\,e_4$ and $U e_3=\ch\,e_3-\sh\,e_4$,
$Ue_0=e_0$, so
\begin{align}
  \langle q_L|U_{\mathbf 5}|t_R\rangle
  &=\big\langle\tfrac{e_0+e_3}{\sqrt2}\,\big|\,\sh\,e_3+\ch\,e_4\big\rangle=\frac{\sh}{\sqrt2},
  &\big|\langle q_L|U|t_R\rangle\big|^2&=\frac{\sh^2}{2},\\
  \langle q_L|U_{\mathbf 5}|q_L\rangle
  &=\tfrac12\big(1+\ch\big)=\cos^2\!\tfrac{h}{2f}.
\end{align}
These are exactly the $\Pi_L\sim\sh^2/2$, $\cos^2(h/2f)$ structures of the thesis $5$-$5$-$5$ form
factors, here \emph{derived} from the overlaps.
\impl{core.overlap\_sym, mchm5.formfactor\_pieces}{test\_coeffs\_5\_5\_5}

\subsection{The $\mathbf{14}$ channel weights and the $4/5$ enhancement}
\label{sec:45}
The $t_R$ embeds in the $\SO4$ singlet of the $\mathbf{14}$, the unit-normalised
symmetric-traceless tensor
\begin{equation}
  S=\frac{1}{\sqrt{20}}\,\mathrm{diag}(1,1,1,1,-4),\qquad \Tr S^2=\frac{4+16}{20}=1 .
\end{equation}
Dress it, $S\mapsto U S U^{\!\top}$, and project onto the $\SO4$ channels of the $\mathbf{14}$. The
squared lengths in the singlet $(1,1)$, bidoublet $(\tfrac12,\tfrac12)$, and $(1,1)\!\to\!(3,3)$
channels are the closed trigonometric \emph{channel weights}
\begin{equation}
  W_{11}=\frac{(4\ch^2-\sh^2)^2}{16},\quad
  W_{22}=\frac{5}{2}\,\sh^2\ch^2,\quad
  W_{33}=\frac{15}{16}\,\sh^4,
\end{equation}
which satisfy unitarity exactly:
\begin{equation}
  W_{11}+W_{22}+W_{33}
  =\Big(1-\tfrac52\sh^2+\tfrac{25}{16}\sh^4\Big)+\tfrac52\sh^2(1-\sh^2)+\tfrac{15}{16}\sh^4=1 .
\end{equation}
The thesis writes the $t_R$ self-energy with an overall coupling $Y_T\sqrt{4/5}$. That
normalisation is \emph{not} an external input: it is the squared index-$4$ component of the
embedding itself,
\begin{equation}
  \boxed{\;\big|S_{44}\big|^2=\Big(\tfrac{-4}{\sqrt{20}}\Big)^2=\frac{16}{20}=\frac45\;}
  \qquad\Longleftrightarrow\qquad \sqrt{\tfrac45}=|S_{44}| .
  \label{eq:45}
\end{equation}
Index $4$ is the coset-singlet direction through which the top mass is generated; the
$\mathbf 5$-singlet has weight $1$ there (the unenhanced reference), so the $4/5$ is a genuine
representation-dependent enhancement. The thesis singlet-channel prefactor then follows as a
\emph{prediction},
\begin{equation}
  \frac{(4\ch^2-\sh^2)^2}{20}\;=\;|S_{44}|^2\,W_{11}\;=\;\frac45\cdot\frac{(4\ch^2-\sh^2)^2}{16},
\end{equation}
which would fail for any inconsistent thesis number --- the check is no longer a tautology that
reads $4/5$ off the prefactor $1/20$.
\impl{decompose.channel\_weights\_sym, core.embedding\_sym}{test\_14\_channel\_weights\_derive\_thesis\_prefactors, test\_4\_5\_is\_derived\_from\_the\_embedding}

\section{Next-to-Minimal Composite Higgs: \texorpdfstring{$\SO6/\SO5$}{SO(6)/SO(5)}}
\label{sec:nmchm}

\subsection{Coset and broken generators}
$\SO6$ acts on indices $0,\dots,5$. The unbroken $\SO5$ acts on $0,\dots,4$; the coset direction is
index $5$. The five broken generators split, under the $\SO4\subset\SO5$, into the Higgs bidoublet
and a singlet:
\begin{equation}
  X_B^{a}=T^{a5}\ (a=0,1,2,3)\quad(\text{the doublet }h_1,\dots,h_4),
  \qquad X_S=T^{45}\quad(\text{the singlet }s).
\end{equation}
This is the thesis decomposition (its 6th index is our index $5$). The $15$ generators close the
$\so6$ algebra; the ten $T^{ab}$ with $a<b\le 4$ are the unbroken $\SO5$.
\impl{so6.gen\_vector, so6.UNBROKEN/BROKEN}{test\_so6\_generators\_close\_the\_algebra, test\_coset\_split\_10\_plus\_5}

\subsection{The five-pNGB Goldstone matrix (closed form)}
In unitary gauge the Higgs points along $\hat a=3$ and the singlet along $X_S$, so the vector
Goldstone is generated by
\begin{equation}
  G=\theta_h\,L^{(3,5)}+\theta_s\,L^{(4,5)},\qquad
  \theta_h=\tfrac{h}{f},\ \ \theta_s=\tfrac{s}{f},
\end{equation}
with $L^{(ab)}$ the real antisymmetric rotation generators. $G$ is supported on the three indices
$\{3,4,5\}$ and has rank $2$, hence satisfies $G^3=-\Theta^2 G$ with
$\Theta=\sqrt{\theta_h^2+\theta_s^2}$. Rodrigues' formula then gives the Goldstone in closed form,
\begin{equation}
  U_6(\theta_h,\theta_s)=\mathbb 1+\frac{\sin\Theta}{\Theta}\,G+\frac{1-\cos\Theta}{\Theta^2}\,G^2 .
  \label{eq:U6}
\end{equation}
Acting on the vacuum $e_5$ reproduces the thesis Goldstone field (eq.~474 there):
\begin{equation}
  \Phi=U_6\,e_5=\Big(0,0,0,\ \theta_h\tfrac{\sin\Theta}{\Theta},\ \theta_s\tfrac{\sin\Theta}{\Theta},\ \cos\Theta\Big)^{\!\top},
\end{equation}
i.e.\ $\Phi=\frac{1}{\Theta}\sin\Theta\,(h_1,\dots,h_4,s,\Theta\cot\Theta)$ with
$\Theta=\varphi/f$, $\varphi=\sqrt{\sum h_i^2+s^2}$. Setting $\theta_s=0$ returns the MCHM planar
rotation \eqref{eq:U5} (with coset index $4\to5$).
\impl{so6.U6\_vector, so6.U6\_vector\_sym, so6.goldstone\_vacuum}{test\_goldstone\_vacuum\_matches\_thesis\_eq474, test\_U6\_vector\_orthogonal\_identity\_homomorphism}

\subsection{The representation tower}
The fermion-partner irreps of $\SO6\cong\SU4$ are built and dressed by the same machinery:
tensors via $n=6$, spinors via the Clifford algebra. Their dimensions, $\SO5$ branchings (from the
quadratic Casimir $C_5=\sum_{a<b\le4}(T^{ab})^2$), and $\SO4$ content are
\begin{center}
\begin{tabular}{llll}
\toprule
irrep & construction & $\dim$ & $\SO5$ branching\\
\midrule
$\mathbf 6$ & vector & $6$ & $\mathbf 5\oplus\mathbf 1$\\
$\mathbf{15}$ & antisym.\ rank 2 (adjoint) & $15$ & $\mathbf{10}\oplus\mathbf 5$\\
$\mathbf{20'}$ & sym.-traceless rank 2 & $20$ & $\mathbf{14}\oplus\mathbf 5\oplus\mathbf 1$\\
$\mathbf{10},\overline{\mathbf{10}}$ & (anti)self-dual 3-form & $10$ & $\mathbf{10}$ (each)\\
$\mathbf 4,\overline{\mathbf 4}$ & Weyl spinor & $4$ & $\mathbf 4$\\
\bottomrule
\end{tabular}
\end{center}
The Casimir calibration $C_5\in\{0,\tfrac52,4,6,10\}\leftrightarrow\{\mathbf1,\mathbf4,\mathbf5,\mathbf{10},\mathbf{14}\}$
identifies each sub-irrep with its multiplicity; the antisymmetric 3-form is $20$-dimensional and
splits, by the $\pm i$ eigenspaces of the Hodge dual $(\star T)_{ijk}=\tfrac{1}{3!}\epsilon_{ijklmn}T_{lmn}$,
into the self-dual $\mathbf{10}$ and anti-self-dual $\overline{\mathbf{10}}$ (the $\SU4$
$\mathbf{10}+\overline{\mathbf{10}}$). The spinors come from six $8\times8$ Euclidean gammas
$\Gamma^A=(\sigma_1\!\otimes\!\gamma^a_{\SO5},\,\sigma_2\!\otimes\!\mathbb 1)$, split by chirality
$\chi=\sigma_3\otimes\mathbb 1$; the $\mathbf 4$ carries no $(\tfrac12,\tfrac12)$ bidoublet, which
is exactly the thesis reason the NM4DCHM uses the $\mathbf 6$ (not the $\mathbf 4$) for $q_L$.
\impl{so6.rep\_basis, so6.U6\_rep, so6.so5\_content, so6\_spinors.py}{test\_so6.py, test\_so6\_spinors.py}

\subsection{Two-field dressing and the singlet pNGB}
The $\mathbf 6$ Goldstone dresses the elementary--composite overlaps with \emph{both} angles. With
$E_{q_L}=(e_0+e_3)/\sqrt2$ and $E_{t_R}=e_5$, no embedding has support on index $4$, so every
overlap depends on $\theta_h,\theta_s$ only through quantities even in $\theta_s$; the potential is
even in $s$, the vacuum sits at $\langle s\rangle=0$, and the singlet mass is the $s$-curvature of
the Coleman--Weinberg potential there. Because the $q_L$ bidoublet dressing carries explicit
$\theta_h$-dependence distinct from the radial $\Theta$ the $t_R$ feels, $V(h,s)$ is genuinely
two-dimensional and the singlet acquires a finite, calculable mass --- the defining NMCHM
observable, absent in the MCHM. At $\theta_s=0$ the mass matrices reduce \emph{entry-for-entry} to
the $\code{pypngb}$-anchored $5$-$5$-$5$ model.
\impl{nmchm6.py (spec, assemble, potential, singlet\_mass2)}{test\_nmchm6.py (reduction, even-in-$s$, singlet mass)}

\section{The Coleman--Weinberg potential and observables}
The one-loop pNGB potential is the supertraced kernel
\begin{equation}
  V(\sh)=\sum_i \frac{c_i}{64\pi^2}\,m_i^2(\sh)^2\log m_i^2(\sh),
  \qquad c_i=\{+3,+6,-12\}
\end{equation}
for $\{$neutral vector, charged vector, coloured Dirac$\}$, with $m_i(\sh)$ the eigenvalues of the
$\sh$-dependent mass matrices. Writing $V=-\gamma\,\sh^2+\beta\,\sh^4$, the vacuum, Higgs mass and
electroweak scale are the closed forms
\begin{equation}
  \xi=\sin^2\!\frac{\langle h\rangle}{f}=\frac{\gamma}{2\beta},\qquad
  m_h^2=\frac{8\beta}{f^2}\,\xi(1-\xi),\qquad v_{\mathrm{EW}}=f\sqrt{\xi},
\end{equation}
and the Barbieri--Giudice tuning is $\Delta_{\rm BG}=\max_i|\partial\ln f/\partial\ln x_i|$ over the
fundamental Lagrangian masses $x_i$. These relations are validated symbolically; the full mass
matrices and the resulting $\xi,m_t,m_h,\Delta_{\rm BG}$ are anchored numerically to the
independent \code{pypngb} engine to $<0.1\%$.
\impl{routes.py, potential.py, spectrum.py, tuning.py}{test\_cw\_kernel\_and\_coefficients, test\_higgs\_mass\_scaling\_closed\_form, test\_anchors.py}

\section{The form-factor route and the oracle boundary}
\label{sec:ff}
The thesis App.~A7 form factors $A_L,A_R,A_M,B$ are transcribed in \code{mchm5}; transcription
proves faithful copying, not physical correctness. They are confronted with the
\code{pypngb}-anchored eigenvalue route by comparing the top mass two ways. With the
\emph{custodial} parametrisation the form-factor route requires (a single $q_L$ coupling $L_q$,
i.e.\ $\Delta_{uL}=\Delta_{dL}$), the two coincide as $\sh\to0$,
\begin{equation}
  \frac{m_t^{\rm ff}(\sh)}{m_t^{\rm eig}(\sh)}=1-0.755\,\sh^2+\mathcal O(\sh^4)\ \xrightarrow[\sh\to0]{}\ 0.999997,
\end{equation}
the finite-$\sh$ gap being the form-factor route's leading-order-in-$\sh^2$ truncation of the exact
(non-polynomial) eigenvalue mass. The App.~A7 form factors are therefore the correct
leading-order two-point functions of the \code{pypngb}-anchored mass matrix --- validated, not
merely transcribed.
\impl{mchm5.formfactor\_pieces, mchm5.fermion\_mass}{test\_formfactor\_route.py}

\section*{Provenance: derived vs.\ input vs.\ anchored}
\addcontentsline{toc}{section}{Provenance}
\begin{description}[leftmargin=1.4em,style=nextline]
  \item[Derived (this document)] the Goldstone matrices $U,U_6$; the rep lifts; the $\SO4/\SO5$
    branchings; the $5$-$5$-$5$ coefficients; the $\mathbf{14}$ channel weights and the
    $4/5=|S_{44}|^2$ enhancement; the full $\SO6$ tower. Each closes symbolically (a named test
    asserts \texttt{simplify}$(\cdot)=0$).
  \item[Input (model structure)] the partner content --- which composite multiplets exist, the
    embeddings, the partner masses. Taken from the thesis/\code{pypngb}, not derived.
  \item[Anchored] the full mass matrices and the physical observables, validated against
    \code{pypngb} ($<0.1\%$). The NMCHM has no public engine and is anchored indirectly by its
    exact reduction to the $5$-$5$-$5$ at $\langle s\rangle=0$, with the new singlet observable
    cross-checked internally (stencil independence + parabolic fit).
\end{description}
After \S\ref{sec:45} and \S\ref{sec:ff}, no quantity that feeds an observable is a thesis
read-off: correctness rests on \code{pypngb} $+$ group theory.

\end{document}
